\documentclass[12pt]{article}
\usepackage[utf8]{inputenc}
\usepackage[T1]{fontenc}
\usepackage{hyperref}
\title{Gugu Warra Language Resources for K-12 Curriculum Development}
\author{Compiled by Claude Sonnet 5.5 (claude-sonnet-5-5) for Andr\'as Kornai}
\date{}
\begin{document}
\maketitle

\section{Introduction}
Gugu Warra (ISO 639-3 wrw; Glottocode gugu1256) is an extinct Pama-Nyungan language of Australia, associated with Cape York Peninsula in northern Queensland. The name (Kuku-Warra, Kuku-Wara) means roughly ``unintelligible speech'', as opposed to Gugu Mini, ``intelligible speech''. It has been applied to several different varieties. The sources I found disagree on the coding. One lists Mbariman-Gudhinma (zmv) as the language of a related Wikipedia page, and Glottolog lists a separate entry for Rimanggudinhma (mbar1253). The Lamalama people are linked with some of these dialects. The language is extinct, so this is a heritage and revitalization case that depends on archival material.

\section{Primary Educational Resources}
The searches found no primers, textbooks or school programmes for Gugu Warra.

\section{Linguistic and Reference Materials}
\begin{itemize}
\item Wikipedia, Mbariman-Gudhinma language, a related variety sometimes labelled Gugu Warra: \url{https://en.wikipedia.org/wiki/Mbariman-Gudhinma_language}
\item University of Sydney archive catalogue entries turned up in the search. I did not open them or confirm what they hold. Example: \url{https://archives-search.sydney.edu.au/nodes/view/17326}
\end{itemize}

\section{Religious and Community Resources}
No scripture or community organisation was identified in the searches.

\section{Digital and Audio Resources}
No audio or digital resources were found. The Sydney archive entries may hold recordings or notes, but I did not check them.

\section{Curriculum Development Context}
Little is documented, and the label covers a loosely defined group of varieties. Any revitalization work should begin by working with the descendant communities (for example Lamalama people). It should also use archival searches, including the Sydney archives, to establish which variety is meant and what material survives. I would not present this as a language with teaching resources.
\end{document}